\uuid{Eh26}
\exo7id{3879}
\titre{exo7 3879}
\auteur{quercia}
\organisation{exo7}
\datecreate{2010-03-11}
\isIndication{false}
\isCorrection{false}
\chapitre{Continuité, limite et étude de fonctions réelles}
\sousChapitre{Etude de fonctions}
\cpge{mpsi-mp2i.fonctions-valeurs-reelles-complexes.a, pcsi.techniques-calcul-fonctions.a, ptsi.techniques-calcul-fonctions.a, mpsi-mp2i.limites-continuite.c, pcsi.limites-continuite.c, ptsi.limites-continuite.c, bcpst1.fonctions-reelles-usuelles.a, bcpst1.limites-continuite.c}
\module{Analyse}
\niveau{L1}
\difficulte{}

\contenu{
\question{
Soit $f : {\R} \to {\R.}$
On dit que $f$ est localement croissante si :
$\forall\ x \in \R,\ \exists\ \varepsilon > 0 \text{ tq }
f_{|]x-\varepsilon, x+\varepsilon[}$ est croissante.

Montrer que : ($f$ est localement croissante) $ \Rightarrow $ ($f$ est croissante).

(\'Etudier $E = \{ x \ge 0 \text{ tq } f_{|[0,x]}\text{ est croissante}\}$ et
         $F = \{ x \le 0 \text{ tq } f_{|[x,0]}\text{ est croissante}\}$)
}
}
